\section{Konzept}

\begin{itemize}
    \item Windows erstellt automatisch Schattenkopien vor Updates und Neustarts
    \item Byte-für-Byte Kopie des Volumes zu einem bestimmten Zeitpunkt
    \item \textbf{Windows 7:} Wiederherstellung per CD
    \item \textbf{Windows 10:} Standardmäßig verfügbar
    \item \textbf{Windows 11:} Standardmäßig deaktiviert
    \item Forensisch wertvoll: gelöschte Dateien können in älteren Schattenkopien noch existieren
\end{itemize}

\section{Workflow}

\begin{lstlisting}[language=bash]
# Shadow Copies auflisten
vshadowinfo /dev/loopXpN
# Ausgabe: Anzahl Stores, Timestamps, GUIDs

# Alle Stores mounten
mkdir -p /tmp/vshadow
vshadowmount /dev/loopXpN /tmp/vshadow
ls /tmp/vshadow/
# -> vss1, vss2, vss3, ... (jeder Store = ein Zeitpunkt)

# Inhalte durchsuchen
fls -pr /tmp/vshadow/vss1 | grep -i 'gesuchte_datei'

# Datei aus Shadow Copy extrahieren
icat /tmp/vshadow/vss3 <INODE> > /tmp/recovered_file
\end{lstlisting}

\section{Anwendungsfall: PST-Recovery}

Gelöschte Outlook-PST-Dateien über Shadow Copies wiederherstellen:

\begin{lstlisting}[language=bash]
# PST in allen Shadow Copies suchen
for i in /tmp/vshadow/vss*; do
  echo "=== $i ==="
  fls -pr "$i" | grep -i '\.pst'
done

# PST extrahieren
icat /tmp/vshadow/vssN <INODE> > /tmp/recovered.pst

# PST analysieren
readpst -o /tmp/pst_out /tmp/recovered.pst
cat /tmp/pst_out/*.mbox | less
\end{lstlisting}

Alternativer Viewer mit mehr Details:

\begin{lstlisting}[language=bash]
pffinfo /tmp/recovered.pst
pffexport /tmp/recovered.pst
\end{lstlisting}

\section{Aufräumen}

\begin{lstlisting}[language=bash]
umount /tmp/vshadow
\end{lstlisting}
